% GENERATED FILE. Edit canonical lessons, then run npm run generate:books.
% canonical-source-hash: fnv1a64:f574f805ec2bc52f
% canonical-lessons: IT-C308-battesimo, IT-C308-battezzare, IT-C308-padrino, IT-C308-madrina, IT-C308-cresima

\chapter{Il battesimo}
\label{ch:it-il-battesimo}
\hlchaptermodality{308}

\begin{chapteropening}
\textbf{By the end of this chapter:} \emph{I can talk about a christening: christening a child, the godparents and confirmation.}

Complete the last lesson of chapter 308: I can talk about a christening: christening a child, the godparents and confirmation.
\end{chapteropening}

\section[il battesimo]{il battesimo — the christening, the baptism}

\label{lesson:IT-C308-battesimo}



\begin{quote}
Before the new one: say the Italian for to wean, then the Italian for the dummy.
\end{quote}

\subsection*{You'll want to know: il battesimo}
\textbf{il battesimo} — \textquotedblleft{}the christening, the baptism\textquotedblright{}.

\begin{quote}
Il battesimo è domenica in chiesa.
\end{quote}

\begin{quote}
\emph{The christening is on Sunday in church.}
\end{quote}

\emph{Il battesimo del volo} — your first flight.

\begin{cousinweb}
\emph{battesimo} comes from Church Latin \emph{baptismus}, from Greek \emph{baptismós} \textquotedblleft{}a dipping\textquotedblright{} (Lewis \& Short).
\end{cousinweb}

\begin{grammarlens}[title={What you've built}]
One more word for telling, explaining and describing at B1.
\end{grammarlens}

\subsection*{Your turn}
\begin{itemize}
\raggedright
  \item \emph{Say it:} \emph{il battesimo}
  \item \emph{Say it:} \emph{il battesimo}, once more
  \item \emph{Say it:} a sentence of your own with \emph{il battesimo}
  \item \emph{Recall:} say the Italian for to wean, then the Italian for the dummy, then say \emph{il battesimo} again
\end{itemize}

\begin{tcolorbox}[breakable,colback=teal!4,colframe=teal!35!black,title={Before you move on}]
What does il battesimo mean? (\textquotedblleft{}The christening, the baptism\textquotedblright{}.) Say it once more.
\end{tcolorbox}
\section[battezzare]{battezzare — to christen, to baptise}

\label{lesson:IT-C308-battezzare}



\begin{quote}
Before the new one: say the Italian for the dummy, then the Italian for the christening.
\end{quote}

\subsection*{You'll want to know: battezzare}
\textbf{battezzare} — \textquotedblleft{}to christen, to baptise\textquotedblright{}.

\begin{quote}
Lo hanno battezzato a giugno.
\end{quote}

\begin{quote}
\emph{They christened him in June.}
\end{quote}

Also to give a name: \emph{hanno battezzato la nave}.

\begin{cousinweb}
\emph{battezzare} is inherited from Church Latin \emph{baptizāre}, from Greek \emph{baptízein} \textquotedblleft{}to dip\textquotedblright{} (Lewis \& Short).
\end{cousinweb}

\begin{grammarlens}[title={What you've built}]
One more word for telling, explaining and describing at B1.
\end{grammarlens}

\subsection*{Your turn}
\begin{itemize}
\raggedright
  \item \emph{Say it:} \emph{battezzare}
  \item \emph{Say it:} \emph{battezzare}, once more
  \item \emph{Say it:} a sentence of your own with \emph{battezzare}
  \item \emph{Recall:} say the Italian for the dummy, then the Italian for the christening, then say \emph{battezzare} again
\end{itemize}

\begin{tcolorbox}[breakable,colback=teal!4,colframe=teal!35!black,title={Before you move on}]
What does battezzare mean? (\textquotedblleft{}To christen, to baptise\textquotedblright{}.) Say it once more.
\end{tcolorbox}
\section[il padrino]{il padrino — the godfather}

\label{lesson:IT-C308-padrino}



\begin{quote}
Before the new one: say the Italian for the christening, then the Italian for to christen.
\end{quote}

\subsection*{You'll want to know: il padrino}
\textbf{il padrino} — \textquotedblleft{}the godfather\textquotedblright{}.

\begin{quote}
Mio zio è il mio padrino.
\end{quote}

\begin{quote}
\emph{My uncle is my godfather.}
\end{quote}

A little \emph{padre}, ‘father’.

\begin{cousinweb}
\emph{padrino} is inherited from Vulgar Latin \emph{patrīnus} (EtymDB), built on \emph{pater} \textquotedblleft{}father\textquotedblright{}.
\end{cousinweb}

\begin{grammarlens}[title={What you've built}]
One more word for telling, explaining and describing at B1.
\end{grammarlens}

\subsection*{Your turn}
\begin{itemize}
\raggedright
  \item \emph{Say it:} \emph{il padrino}
  \item \emph{Say it:} \emph{il padrino}, once more
  \item \emph{Say it:} a sentence of your own with \emph{il padrino}
  \item \emph{Recall:} say the Italian for the christening, then the Italian for to christen, then say \emph{il padrino} again
\end{itemize}

\begin{tcolorbox}[breakable,colback=teal!4,colframe=teal!35!black,title={Before you move on}]
What does il padrino mean? (\textquotedblleft{}The godfather\textquotedblright{}.) Say it once more.
\end{tcolorbox}
\section[la madrina]{la madrina — the godmother}

\label{lesson:IT-C308-madrina}



\begin{quote}
Before the new one: say the Italian for to christen, then the Italian for the godfather.
\end{quote}

\subsection*{You'll want to know: la madrina}
\textbf{la madrina} — \textquotedblleft{}the godmother\textquotedblright{}.

\begin{quote}
La madrina mi ha regalato una catenina d'oro.
\end{quote}

\begin{quote}
\emph{My godmother gave me a little gold chain.}
\end{quote}

A new ship has \emph{una madrina} too, who names it.

\begin{cousinweb}
\emph{madrina} is the feminine pair of \emph{padrino}, built on \emph{madre} \textquotedblleft{}mother\textquotedblright{}, from Latin \emph{māter} (Lewis \& Short).
\end{cousinweb}

\begin{grammarlens}[title={What you've built}]
One more word for telling, explaining and describing at B1.
\end{grammarlens}

\subsection*{Your turn}
\begin{itemize}
\raggedright
  \item \emph{Say it:} \emph{la madrina}
  \item \emph{Say it:} \emph{la madrina}, once more
  \item \emph{Say it:} a sentence of your own with \emph{la madrina}
  \item \emph{Recall:} say the Italian for to christen, then the Italian for the godfather, then say \emph{la madrina} again
\end{itemize}

\begin{tcolorbox}[breakable,colback=teal!4,colframe=teal!35!black,title={Before you move on}]
What does la madrina mean? (\textquotedblleft{}The godmother\textquotedblright{}.) Say it once more.
\end{tcolorbox}
\section[la cresima]{la cresima — the confirmation (sacrament)}

\label{lesson:IT-C308-cresima}



\begin{quote}
Before the new one: say the Italian for the godfather, then the Italian for the godmother.
\end{quote}

\subsection*{You'll want to know: la cresima}
\textbf{la cresima} — \textquotedblleft{}the confirmation (sacrament)\textquotedblright{}.

\begin{quote}
A dodici anni ho fatto la cresima.
\end{quote}

\begin{quote}
\emph{I was confirmed at twelve.}
\end{quote}

\emph{Fare la cresima} — to be confirmed.

\begin{cousinweb}
\emph{cresima} comes from Church Latin \emph{chrisma} \textquotedblleft{}anointing oil\textquotedblright{}, from Greek \emph{khrîsma} \textquotedblleft{}an ointment\textquotedblright{} (Lewis \& Short).
\end{cousinweb}

\begin{grammarlens}[title={What you've built}]
One more word for telling, explaining and describing at B1.
\end{grammarlens}

\subsection*{Your turn}
\begin{itemize}
\raggedright
  \item \emph{Say it:} \emph{la cresima}
  \item \emph{Say it:} \emph{la cresima}, once more
  \item \emph{Say it:} a sentence of your own with \emph{la cresima}
  \item \emph{Recall:} say the Italian for the godfather, then the Italian for the godmother, then say \emph{la cresima} again
\end{itemize}

\begin{tcolorbox}[breakable,colback=teal!4,colframe=teal!35!black,title={Before you move on}]
What does la cresima mean? (\textquotedblleft{}The confirmation (sacrament)\textquotedblright{}.) Say it once more.
\end{tcolorbox}
